\documentclass[11pt]{article}

% ============================================================
% Packages
% ============================================================
\usepackage[margin=1in]{geometry}
\usepackage{graphicx}
\usepackage{booktabs}
\usepackage{amsmath}
\usepackage{amssymb}
\usepackage{siunitx}
\usepackage{float}
\usepackage{hyperref}
\usepackage{xcolor}
\usepackage{listings}
\usepackage{enumitem}
\usepackage{caption}
\usepackage{subcaption}

% ============================================================
% Hyperlinks
% ============================================================
\hypersetup{
    colorlinks=true,
    linkcolor=blue,
    urlcolor=blue,
    citecolor=blue
}

% ============================================================
% Code formatting
% ============================================================
\lstset{
    basicstyle=\ttfamily\small,
    frame=single,
    breaklines=true,
    columns=fullflexible,
    showstringspaces=false
}

% ============================================================
% Document
% ============================================================
\begin{document}

\begin{center}
    {\LARGE \textbf{ECSE 4320 Advanced Computer Systems}}\\[0.4em]
    {\Large Project \#1: SIMD Advantage Profiling}\\[1em]

    \textbf{Your Name}\\
    Rensselaer Polytechnic Institute\\
    Fall 2026\\[0.5em]

    \textbf{GitHub Repository:}
    \url{https://github.com/YOUR-USERNAME/YOUR-REPOSITORY}
\end{center}

\vspace{1em}

% ============================================================
\section{Introduction}
% ============================================================

% Briefly explain the purpose of the project.
%
% Questions to answer:
% - What is SIMD?
% - What are you trying to measure?
% - Why might SIMD speedup differ from the theoretical vector width?
% - What major factors will you investigate?
%
% Keep this concise. This is not meant to be a long literature review.

[Write introduction here.]

% ============================================================
\section{Experimental Platform}
% ============================================================

All experiments were performed on the system described in
Table~\ref{tab:system}.

\begin{table}[H]
    \centering
    \caption{Experimental system configuration.}
    \label{tab:system}
    \begin{tabular}{ll}
        \toprule
        \textbf{Component} & \textbf{Configuration} \\
        \midrule
        CPU model & AMD Ryzen 7 260 \\
        Physical cores & 8 \\
        Threads & 16 \\
        Logical CPUs / SMT & 16 \\
        SIMD ISA support & SSE, AVX, AVX2, AVX-512 \\
        L1 cache & 32 KB per core \\
        L2 cache & 1 MB per core \\
        LLC / L3 cache & 16 MiB shared \\
        RAM & 30 GB \\
        Operating system & Linux Mint 22.3 \\
        Kernel version & 7.0.0-38-generic \\
        Compiler & GCC 13.3.0 \\
        Python & 3.12.9 \\
        CPU frequency policy/governor & powersave \\
        \bottomrule
    \end{tabular}
\end{table}

% Explain HOW you obtained this information.
% Do not simply fill the table with values you cannot explain.

% ============================================================
\section{Benchmark Methodology}
% ============================================================

\subsection{Selected Kernels}

The following three or more computational kernels were evaluated:

\begin{enumerate}
    \item [Kernel 1]
    \item [Kernel 2]
    \item [Kernel 3]
\end{enumerate}

% Briefly define each kernel mathematically or with concise pseudocode.
%
% The project permits:
% - SAXPY / AXPY
% - dot product / reduction
% - elementwise multiplication
% - 1D three-point stencil

\subsection{Correctness Validation}

% Explain how you established that optimized/vectorized results
% were numerically correct.
%
% Questions:
% - What reference implementation did you use?
% - How were results compared?
% - Did floating-point rounding require a tolerance?
% - Were multiple/random inputs tested?

[Describe correctness procedure.]

\subsection{Timing Method}

% Explain exactly what is timed.
%
% Include:
% - timer used
% - whether initialization is outside timed region
% - warm-up procedure, if any
% - number of repetitions
% - statistic reported (median, mean, minimum, etc.)
% - how measurement variability is represented

[Describe timing methodology.]

\subsection{Compiler Configurations}

Table~\ref{tab:compiler} summarizes the compiler configurations used
for comparison.

\begin{table}[H]
    \centering
    \caption{Compiler configurations.}
    \label{tab:compiler}
    \begin{tabular}{lll}
        \toprule
        \textbf{Configuration} & \textbf{Compiler Flags} & \textbf{Purpose} \\
        \midrule
        Scalar-only & [Fill in] & Disable vectorization \\
        Auto-vectorized & [Fill in] & Allow compiler vectorization \\
        CPU-targeted & [Fill in] & Target available CPU features \\
        \bottomrule
    \end{tabular}
\end{table}

% Do not copy flags here until you understand what each one does.
% Record the exact command lines used in the experiment.

% ============================================================
\section{Experiment 1: Scalar vs.\ SIMD}
% ============================================================

\subsection{Question}

% What specific question does this experiment answer?

[Write experimental question.]

\subsection{Hypothesis}

% IMPORTANT: Write what you expected BEFORE discussing the results.
% Explain WHY based on architecture.

[Write pre-experiment hypothesis.]

\subsection{Experimental Setup}

% State:
% - which kernels
% - N / working-set sizes
% - datatype
% - compiler configurations
% - controlled variables
% - repetitions
% - metric(s)

[Describe setup.]

\subsection{Results}

% Replace placeholder filename with your generated plot.

The SIMD speedup is calculated as

\begin{equation}
    S_{\mathrm{SIMD}}
    =
    \frac{T_{\mathrm{scalar}}}
         {T_{\mathrm{SIMD}}}.
\end{equation}

\subsection{Analysis}

% Don't merely say "SIMD was faster."
%
% Explain:
% - How large was the speedup?
% - Did it change with N?
% - Did different kernels behave differently?
% - Was speedup close to vector width?
% - What architectural limitation(s) explain the result?

[Analyze results.]

% ============================================================
\section{Experiment 2: Working-Set Size and Cache Locality}
% ============================================================

\subsection{Question}
[Write question.]

\subsection{Hypothesis}
[Write hypothesis before discussing results.]

\subsection{Experimental Setup}

% Sweep N from small arrays through increasingly large working sets.
% Explain how the chosen sizes relate to L1, L2, LLC, and DRAM.

[Describe setup.]

\subsection{Results}
x

\subsection{Analysis}
x
% Discuss:
% - transitions associated with cache levels
% - when SIMD gains compress
% - evidence that the workload becomes memory-bound
% - relevant performance-counter evidence, if collected

[Analyze results.]

% ============================================================
\section{Experiment 3: Alignment and Tail Handling}
% ============================================================

\subsection{Question}
[Write question.]

\subsection{Hypothesis}
[Write hypothesis.]

\subsection{Experimental Setup}

% Compare:
% - aligned vs intentionally shifted/misaligned data
% - sizes divisible by vector width vs sizes with a remainder

[Describe setup.]

\subsection{Results}
x

\subsection{Analysis}
x
% Possible things to investigate:
% - unaligned loads/stores
% - scalar cleanup loops
% - masking
% - prologue/epilogue overhead
%
% Only claim mechanisms supported by your evidence.

[Analyze results.]

% ============================================================
\section{Experiment 4: Stride and Access Pattern}
% ============================================================

\subsection{Question}
[Write question.]

\subsection{Hypothesis}
[Write hypothesis.]

\subsection{Experimental Setup}

% Typical comparison:
% stride = 1, 2, 4, 8
%
% State exactly how your benchmark defines stride so that
% comparisons remain fair.

[Describe setup.]

\subsection{Results}
x

\subsection{Analysis}

% Discuss:
% - spatial locality
% - effective memory bandwidth
% - cache-line utilization
% - SIMD efficiency
% - gather-like behavior if applicable

[Analyze results.]

% ============================================================
\section{Experiment 5: Data Type}
% ============================================================

\subsection{Question}
[Write question.]

\subsection{Hypothesis}
[Write hypothesis.]

\subsection{Experimental Setup}

% Required comparison includes float32 vs float64.
% Explain how everything except datatype was controlled.

[Describe setup.]

\subsection{Results}
x

\subsection{Analysis}

% Discuss:
% - SIMD lane count
% - bytes moved per element
% - operation cost
% - arithmetic intensity where relevant

[Analyze results.]

% ============================================================
\section{Vectorization Verification}
% ============================================================

% The assignment explicitly requires evidence that SIMD occurred.
%
% Include concise evidence, NOT pages of assembly.
%
% You might include:
% 1. relevant compiler vectorization report
% 2. a very small disassembly excerpt
%
% Explain what the evidence means.

\subsection{Compiler Vectorization Report}

\begin{lstlisting}
[Insert a SHORT relevant compiler-report excerpt.]
\end{lstlisting}

[Explain what the compiler report demonstrates.]

\subsection{Disassembly Evidence}

\begin{lstlisting}
[Insert only a few relevant assembly instructions.]
\end{lstlisting}

% Explain the instructions you identify.
% Do not paste assembly that you cannot explain.

[Explain why this constitutes evidence of vector execution.]

% ============================================================
\section{Performance-Counter Analysis}
% ============================================================

% Include perf or equivalent counters where they contribute to
% your architectural explanation.
%
% Avoid dumping every available counter into a table.
% Choose counters that help answer your experimental questions.

\begin{table}[H]
    \centering
    \caption{Selected performance-counter measurements.}
    \label{tab:counters}
    \begin{tabular}{lrr}
        \toprule
        \textbf{Metric} & \textbf{Case A} & \textbf{Case B} \\
        \midrule
        Cycles & [Fill in] & [Fill in] \\
        Instructions & [Fill in] & [Fill in] \\
        Cache-related metric & [Fill in] & [Fill in] \\
        Other relevant metric & [Fill in] & [Fill in] \\
        \bottomrule
    \end{tabular}
\end{table}

[Interpret the counters rather than simply listing them.]

% ============================================================
\section{Roofline Analysis}
% ============================================================

\subsection{Arithmetic Intensity}

Arithmetic intensity is defined as

\begin{equation}
    AI =
    \frac{\text{floating-point operations}}
         {\text{bytes transferred}}.
\end{equation}

% Derive AI for at least one selected kernel.
% Be explicit about assumptions regarding memory traffic.

[Show your calculation.]

\subsection{Machine Limits}

The Roofline performance bound can be written as

\begin{equation}
    P_{\mathrm{attainable}}
    \leq
    \min\left(
        P_{\mathrm{peak}},
        AI \times B_{\mathrm{memory}}
    \right),
\end{equation}

where $P_{\mathrm{peak}}$ is the estimated peak floating-point
throughput and $B_{\mathrm{memory}}$ is measured memory bandwidth.

% Explain how you obtained BOTH quantities.
% Do not substitute unexplained internet specifications for measured
% values where the assignment specifically asks for measured bandwidth.

[Describe peak-performance and bandwidth estimates.]

\subsection{Interpretation}

% Is your selected kernel compute-bound or memory-bound?
% Does its measured behavior agree with the Roofline reasoning?
% If not, what other limitations may matter?

[Write Roofline interpretation.]

% Optional figure if you construct an actual Roofline plot:
%
% \begin{figure}[H]
%     \centering
%     \includegraphics[width=0.80\textwidth]{figures/roofline.pdf}
%     \caption{Roofline analysis of the selected kernel.}
%     \label{fig:roofline}
% \end{figure}

% ============================================================
\section{Negative / Failure Case}
% ============================================================

\subsection{Hypothesis}

% Identify a situation in which an apparently useful optimization
% may fail, or an apparently SIMD-friendly loop may not behave as expected.

[Write hypothesis.]

\subsection{Results}

[Present measurement/figure/table.]

\subsection{Explanation}

% The project specifically requires a meaningful negative example.
%
% Explain WHY the expected optimization did not help.
% Verify the mechanism where possible rather than guessing.

[Explain failure case.]

% ============================================================
\section{Experimental Limitations and Anomalies}
% ============================================================

% Discuss:
% - unexplained or unexpected measurements
% - measurement noise
% - hardware/tool limitations
% - counters you could not access
% - assumptions
%
% Do not hide anomalous results.

[Discuss limitations and anomalies.]

% ============================================================
\section{AI Use Statement}
% ============================================================

% The assignment explicitly requires this.
%
% State:
% - which AI tool(s) you used
% - what tasks AI assisted with
% - what code/methodology/results you independently tested
% - anything suggested by AI that you modified or rejected
%
% Be specific and truthful.

[Write AI-use statement.]

% ============================================================
\section{Reproducibility}
% ============================================================

% Briefly explain how someone can reproduce your experiments.
%
% Point to:
% - source code
% - build scripts
% - run scripts
% - raw data
% - processed data
% - plotting scripts
% - README

The complete source code, scripts, raw measurements, processed data,
and plotting code are available in the project repository:

\begin{center}
    \url{https://github.com/YOUR-USERNAME/YOUR-REPOSITORY}
\end{center}

% ============================================================
\section{Conclusion}
% ============================================================

% Answer the big questions:
%
% - When did SIMD provide substantial speedup?
% - When did it provide limited speedup?
% - What architectural factors mattered most?
% - What practical lesson did you learn about SIMD optimization?
%
% Keep this concise and quantitative.

[Write conclusion.]

\end{document}